\section*{Introduction générale}
\addcontentsline{toc}{section}{Introduction générale}
\markboth{Introduction générale}{Introduction générale}

La migration internationale est un fait structurant de la société
sénégalaise. Ancienne, d'abord tournée vers les pays voisins et la France,
elle s'est diversifiée depuis les années 1990 vers l'Italie, l'Espagne et
l'Amérique du Nord, portée par des réseaux familiaux et villageois qui en
abaissent le coût pour ceux qui suivent \parencite{massey1993, ndiaye2016}.
Partir n'y est pas une décision individuelle mais une stratégie de ménage,
la famille finançant le voyage d'un de ses membres et attendant de lui, en
retour, un soutien régulier \parencite{stark1985}. Il en résulte une
diaspora nombreuse et durablement installée à l'étranger, dont la
manifestation économique la plus tangible, pour les ménages d'origine, est
le flux de transferts qu'elle génère. Devenus l'une des premières
ressources extérieures du pays, ces envois sont passés de 233 millions de
dollars en 2000 à 2\,220 millions en 2017, soit de 166 à 1\,242 milliards
de FCFA, et atteignent 12,1\,\% du produit intérieur brut, ce qui place le
pays au quatrième rang des récepteurs d'Afrique subsaharienne
\parencite{ansd2018profil}. À l'échelle des ménages, le phénomène n'a rien
de marginal, une part importante des ménages sénégalais déclarant recevoir
un transfert de l'étranger, pour des montants dont l'ordre de grandeur
avoisine le seuil de pauvreté monétaire annuel par tête.

Cette ressource se distingue des autres flux extérieurs par trois traits :
elle est décentralisée, arrivant directement dans le budget du ménage sans
transiter par une administration~; elle est relativement stable, les
envois étant adossés au revenu du travail de personnes installées à
l'étranger \parencite{maimbo2005}~; elle est enfin, pour partie,
contracyclique du point de vue du ménage receveur, l'envoi s'intensifiant
lorsque celui-ci subit un choc \parencite{gubert2002}. Ces trois traits
font des transferts un candidat sérieux au rôle d'amortisseur des
privations. Or les privations touchant les enfants restent d'une ampleur
considérable : selon l'actualisation de l'ANSD et de l'UNICEF
\parencite{ansd2024}, 50,7\,\% des enfants sénégalais de 0 à 17 ans
subissaient en 2018 des privations simultanées dans au moins quatre des
sept domaines retenus pour apprécier leur bien-être, l'éducation, la
santé, la nutrition, l'eau, l'assainissement, le logement et la
protection. Ces manques échappent largement aux mesures fondées sur le
seul revenu du ménage, dont l'approche monétaire suppose une répartition
interne équitable que rien n'assure \parencite{deaton1997}, et ils
appellent une mesure au niveau de l'enfant lui-même, les privations
retenues étant individuelles par construction et les besoins variant
fortement avec l'âge \parencite{deneubourg2012}. À ces raisons de mesure
s'ajoute un motif de fond, les privations subies pendant l'enfance ayant
des effets durables sur le capital humain et se transmettant d'une
génération à la suivante \parencite{gordon2003}.

La capacité de cette ressource privée massive à réduire effectivement les
privations des enfants n'a pourtant rien d'évident. La destination des
fonds dépend du mobile de l'envoi et n'a aucune raison de coïncider avec
les dimensions que retient l'indice, un transfert destiné au soutien
courant du ménage n'améliorant pas nécessairement l'accès à l'eau ou la
scolarisation. Plusieurs privations supposent en outre des investissements
collectifs qu'un flux monétaire privé ne crée pas, et le coût non
monétaire de la migration, l'absence du parent, peut compenser, voire
annuler, l'apport financier qu'elle permet \parencite{mansuri2006,
davis2016}. Deux traits du cas sénégalais donnent à cette question une
acuité particulière : l'ampleur du phénomène, qui fait que la réponse
engage une fraction substantielle des enfants du pays, et le profil
masculin des migrants, qui superpose dans les mêmes ménages l'apport
financier et l'absence du père.

À cette incertitude sur l'effet s'ajoute une difficulté d'identification.
Recevoir un transfert n'est pas un événement fortuit, mais la conséquence
d'une migration antérieure, elle-même conditionnée par l'accès à un réseau,
la richesse initiale et les relations sociales du ménage
\parencite{massey1993}. Les ménages bénéficiaires diffèrent donc des
autres avant même de recevoir le moindre franc, et sur des
caractéristiques dont certaines ne s'observent pas : comparer directement
les enfants des deux groupes attribuerait aux transferts des écarts qui
leur préexistent. L'enjeu dépasse la connaissance académique, les
transferts étant régulièrement présentés comme un levier de réduction de
la pauvreté : si leur effet sur les privations des enfants s'avérait
limité ou concentré sur certaines dimensions, un ménage bénéficiaire ne
saurait être tenu pour couvert ni écarté des dispositifs de protection
sociale, et l'apport privé ne se substituerait pas à l'investissement
collectif dans les services de base. Deux lacunes ressortent par ailleurs
de la littérature existante : une lacune de mesure, les travaux sur les
transferts raisonnant au niveau du ménage et sur des indicateurs
monétaires ou sectoriels, et une lacune d'identification, les études
disponibles en Afrique de l'Ouest reposant le plus souvent sur des données
transversales qui ne permettent pas de neutraliser les caractéristiques
inobservées et stables des ménages.

Ce mémoire est dès lors structuré par l'interrogation suivante : dans
quelle mesure les transferts de migrants influencent-ils la pauvreté
multidimensionnelle des enfants au Sénégal~? Il vise à évaluer l'impact
d'une exposition durable à ces transferts sur les privations des enfants,
mesurées dans les sept dimensions de leur bien-être, et à établir si cet
impact, loin d'être uniforme, varie selon les dimensions et le montant
reçu. Trois tâches en découlent : construire un indice de pauvreté
multidimensionnelle adapté aux enfants sénégalais, décliné par groupe
d'âge, et en mesurer le niveau à chacune des deux éditions~; estimer
l'impact d'une exposition durable aux transferts sur le nombre de
privations, séparément dans chaque groupe d'âge, en corrigeant la
sélection des ménages bénéficiaires~; analyser l'hétérogénéité de cet
impact selon la dimension et l'indicateur de privation, et selon le
montant des transferts reçus.

Deux hypothèses sont mises à l'épreuve. La première porte sur l'existence
de l'effet et découle du rôle d'amortisseur que la littérature prête aux
transferts~; la seconde porte sur la forme de l'effet, dont rien ne
garantit l'uniformité, la destination des fonds et l'offre locale de
services variant d'un contexte à l'autre.

\begin{itemize}
  \item H1) Une exposition durable aux transferts de migrants réduit
        significativement la pauvreté multidimensionnelle des enfants,
        dans chacun des groupes d'âge.
  \item H2) Cet impact est hétérogène selon les dimensions et les
        indicateurs de privation et selon le montant des transferts
        reçus.
\end{itemize}

Pour répondre à cette question, ce mémoire mobilise les deux éditions de
l'Enquête harmonisée sur les conditions de vie des ménages, celle de
2018-2019 et celle de 2021-2022, dont l'appariement permet de suivre
6\,127 ménages et 17\,786 enfants d'une édition à l'autre. La pauvreté est
mesurée selon l'approche MODA, en sept dimensions et treize indicateurs
déclinés par groupe d'âge. L'effet d'une exposition durable aux transferts
est estimé par appariement sur score de propension, sur le nombre de
privations de la seconde édition, le score incluant le niveau initial de
privations et les caractéristiques prédéterminées du ménage~; l'estimation
est conduite séparément dans chaque groupe d'âge, et sa robustesse est
éprouvée par la définition du traitement, le seuil de privation et les
trois algorithmes d'appariement. La contribution de ce travail est double :
sur le plan de la mesure, il applique au Sénégal une mesure de la pauvreté
prenant l'enfant pour unité d'analyse et agrégeant sept dimensions de son
bien-être, là où la littérature sur les transferts s'en tient au ménage et
au revenu~; sur le plan de l'identification, il exploite la dimension
longitudinale des deux éditions pour conditionner l'appariement sur le
niveau initial de privations de chaque enfant et pour restreindre le
traitement aux bénéficiaires stables, ce qu'aucune étude recensée en
Afrique de l'Ouest ne fait sur cette question.

Ce mémoire s'articule en sept chapitres. Le chapitre 1 pose le cadre
conceptuel, en définissant les transferts de migrants et en situant la
pauvreté multidimensionnelle parmi les approches de la pauvreté. Le
chapitre 2 dresse une revue de la littérature théorique et empirique sur
les transferts de migrants et le bien-être des enfants. Le chapitre 3
présente les données EHCVM~I et II ainsi que le cadre méthodologique, avec
la définition de la variable de traitement, la construction de l'indice
MODA et la stratégie d'identification par appariement sur score de
propension. Le chapitre 4 mesure la pauvreté multidimensionnelle des
enfants aux deux éditions et confronte les privations des enfants de
ménages bénéficiaires et non bénéficiaires. Le chapitre 5 présente les
résultats des estimations. Le chapitre 6 discute la portée des résultats
et les limites de l'étude. Le chapitre 7 conclut et formule des
recommandations de politique économique.
